\sectiontheme{emipurple}
\section{Actions}
\label{sec:actions}

Actions are, in effect, functions: they take an input and return an
output, corresponding to a ``controller'' in traditional API frameworks.
Each action is an endpoint, but --- unlike a typical controller --- an
action can also be called from contexts other than HTTP, for example the
CLI, since \texttt{cliName} turns any action into a \texttt{urfave/cli}
command for free.

\begin{lstlisting}
name: sampleModule
actions:
  - name: getSinglePost
    url: https://jsonplaceholder.typicode.com/posts/:id
    cliName: get-single-post
    method: get
    out:
      fields:
        - {name: userId, type: int64}
        - {name: title,  type: string}
  - name: sampleSse
    url: http://localhost:3000/stream
    method: get
    out:
      fields: [{name: message, type: string}]
  - name: webSocketOrgEcho
    url: "wss://echo.websocket.org/.ws"
    method: reactive
    in:
      fields: [{name: firstName, type: string}]
    out:
      fields: [{name: lastName, type: string}]
\end{lstlisting}

\begin{sidebar}{Action properties}
\small
\begin{tabular}{@{}p{0.22\linewidth}p{0.72\linewidth}@{}}
\textbf{name} & required, camelCase, unique within the module \\
\textbf{cliName} & overrides the CLI command name (defaults to \texttt{name}) \\
\textbf{cliShort} & shorter alternative alongside \texttt{cliName} \\
\textbf{url} & HTTP route; omitted means the action is CLI-only \\
\textbf{method} & \texttt{post}/\texttt{patch}/\texttt{put}/\texttt{get}/\texttt{delete}/\texttt{reactive} \\
\textbf{binaryType} & websocket payload type, text by default \\
\textbf{qs} & type-safe query parameters, for both CLI and HTTP \\
\textbf{description} & used in specs and generated documentation \\
\textbf{in} / \textbf{out} & request/response body definitions \\
\end{tabular}
\end{sidebar}

Only \texttt{name} is strictly required for an action to compile;
\texttt{in}, \texttt{out}, \texttt{method} and \texttt{url} are simply the
properties developers reach for most often. Calling a non-listed method
falls back to HTTP with that non-standard verb. \texttt{reactive} is
Emi's own addition on top of the HTTP-standard methods --- it's what
turns an action into a WebSocket channel instead of a request/response
pair (\S\ref{sec:reactive}), and a plain \texttt{get} whose handler
streams chunked output is how an SSE action like \texttt{sampleSse} above
is declared, with no separate keyword needed.
